% 04-experiments.tex — truth: 本文件（2026-09-10 交割：R22 后全部正文修订直改 tex，原装配源 md 已冻结为历史设计稿，见其头部治理注记）
% 纪律: 4.1 数据集与三层口径协议(20/24 有效监督口径随行) / 4.2 实现细节+统计协议 / 4.3 主实验 E1-E6(每实验先声明 claim 句) /
% 4.4 对比研究(NTU 等价性段=成文块装配); tab2 真实 K9 列空缺如实标注; 内部路径/文件名禁入正文(ROUND2 发现#7)。
\section{Experiments}\label{sec:experiments}

\subsection{Datasets and the three-tier protocol}\label{sec:datasets}

Because no dataset we could identify spans the scenarios our claims require (Section~\ref{ssec:rw-animal}), we evaluate under a three-tier protocol; numbers are never mixed within a comparison, and cross-tier comparisons occur only where a pre-registered protocol specifies them and are labeled as such. The \emph{synthetic} tier serves pipeline shake-down and controlled ablations; the \emph{public-real} tier provides label-free and low-resource evaluation on real animal skeletons; the \emph{real-K9} tier (working-dog training footage) is reserved for the target-domain deployment scenario and is not exercised in this study. Two further evaluation contexts are exercised under the same reporting rules and are never pooled with the three tiers: a \emph{synthetic-offset} variant of the synthetic tier (the same generator with Gaussian joint-coordinate noise, $\sigma{=}0.10$ added to skeleton coordinates; its distribution was fixed by a pre-registered diagnostic before any offset experiment ran) used for the distribution-shift and warm-start studies; and public \emph{benchmarks} outside the tiers---NTU60/NTU120/UCF101 (human actions) and PanAf500 (great apes)---used for the pre-registered E9-series cross-domain retention studies, each reported against its own full-budget reference. The working-dog pilot is nonetheless pre-registered before any data collection (protocol PSD-K9-PREREG-001, committed to the public review repository): it fixes a double-annotation agreement gate (Cohen's $\kappa\geq 0.60$), session-level splits, a single primary endpoint with a paired decision rule, and a no-redistribution clause for raw footage and derived skeletons, so that target-domain evidence, when it arrives, cannot be reported selectively after the fact.

The public-real tier draws on three sources. InterPet4D~\citep{peng2026interpet4d} provides 226 SMAL-fitted skeleton sequences of human--pet interactions (225 used; one all-NaN sequence excluded and listed) and serves as the unlabeled stream for self-supervised pretraining and segmentation evaluation. From the official Animal Kingdom action-recognition lists~\citep{ng2022animalkingdom} we aggregate the canine-species subset: 329 videos (231 train / 98 test) with 34{,}772 frame-level annotation rows, used for weak-supervision and behavior-classification evaluation. APTv2~\citep{yang2023aptv2} supplies a large pose pool (official counts: 2{,}749 clips, 41{,}235 frames; our local inventory holds 41{,}179 images; consumed unlabeled in this study) for expansion experiments. Skeletons are represented as world-coordinate keypoint sequences of shape $T{\times}24{\times}3$ with a per-joint confidence channel. One structural caveat applies to every real-domain skeleton product in this paper: the dog-pose annotations never label four of the 24 graph slots (both eyes, withers, and throat), so real-domain keypoints are topology-isomorphic with 20/24 effective supervision channels; the four dead channels are hard-masked at assembly and NaN-guarded downstream. A pre-registered expansion tier (v2; protocol PSD-AKV2-PREREG-001, frozen before the build) re-extracts the same canine videos as multiple contiguous segments under a per-frame label-consistency gate, yielding 352 clips in an eight-class space; it serves as a pre-registered parallel replication layer and never replaces v1 numbers (experiment E7b).

\subsection{Implementation details}\label{sec:implementation}

All experiments run on a single consumer GPU (RTX 5060 Laptop, 8\,GB); self-supervised pretraining takes about 13 minutes for 120 epochs on InterPet4D, which we report as evidence that the physics layer is reproducible without data-center resources. Pretraining uses a learning rate of 0.05, a queue of 1{,}024, and window length $T{=}64$; this configuration is the first working setting, not the outcome of a systematic search. $k$-NN probes report 5-fold mean$\pm$std; end-to-end arms report seed-level dispersion, which is a spread estimate (sample standard deviation, ddof~1) and not a confidence interval; per-experiment seed choices are listed with their sources in the released experiment ledger and protocol files. Following our statistics protocol, aggregate comparisons are accompanied by paired seed-level tests; within each experiment family the reported $p$ values are Holm--Bonferroni corrected (corrected values printed alongside the test), effect sizes are reported alongside raw values with a random-baseline reference, and outliers are explicitly explained rather than silently removed. The E9-series correction family (Section~\ref{sec:comparison}) consolidates the series' eight paired tests: the pipeline-versus-linear gap tests on three tiers (NTU60, NTU120, UCF101; one-sample $t$ and Wilcoxon signed-rank each) plus the two external-SSL Wilcoxon tests of Section~\ref{sec:comparison}. Each member's decision rule was fixed by its own pre-registered protocol before its run; the consolidated eight-test family was fixed before corrected inference entered this manuscript, and the consolidation is dated in the repository (2026-09-07). The PanAf500 gap is reported directionally and is not a family member (its arm distribution ties with the linear arm, so a signed-rank test is uninformative; raw $p=0.75$), and adding it as a ninth member changes no corrected verdict; under the narrower six-test gap-only sub-family every qualitative conclusion is also unchanged. Registration fixes decision rules before outcomes exist and prevents selective reporting; it did not catch the protocol-implementation errors we corrected (the oracle-label leakage of the initial end-to-end runs, Section~\ref{ssec:proto}), which came from internal adversarial review. Protocol freeze is evidenced by the dated commit history of the public review repository (tag \texttt{review-snapshot}); one protocol (the v2 expansion) and its results landed in the same repository batch, so its freeze is bookkeeping-grade evidence rather than third-party timestamping.

\subsection{Main results}\label{sec:main-results}

Table~\ref{tab:main} consolidates the main results across the three tiers plus the public human-domain benchmark used for implementation-equivalence verification; each numbered experiment below opens by stating the claim it tests; experiment labels follow the released experiment ledger (E5, the active-learning study, is reported as an exploratory finding in Section~\ref{sec:ablation} rather than as a numbered paragraph here; E8, the internal identifier for the budget curve, is folded into the E7 discussion). Claim tags used throughout: \textbf{C1} decoupled transition cost, \textbf{C3} pretraining value, \textbf{C4} unsupervised segmentation, \textbf{C5} anchor-guided semantic expansion, \textbf{C6} low-budget self-training behavior---approaching full supervision on the synthetic-offset and human tiers, near-chance on the canine tier (tier-dependent; see E7), \textbf{C7} active learning (exploratory negative); \textbf{C2} is the novelty claim, tested by the survey (Appendix~\ref{app:novelty}) rather than by an experiment.

\paragraph{E1: representation quality of the physics layer (C3).}
This experiment tests whether self-supervised pretraining on unlabeled quadruped skeletons yields discriminative dynamics representations. Using subject identity as a $k$-NN probe on InterPet4D---a proxy task that measures representation discriminability, not behavior accuracy---pretraining reaches \textbf{20.89\% $\pm$ 4.04\%} top-1 (sample std over the five folds) against an 8.33\% random baseline ($2.51\times$). The fold-level range is wide (15.56--26.67\%), attributed to the small number of sequences and subject diversity; a projection-head feature caliber yields 12.89\% ($1.55\times$), recorded in the released run report, and both calibers are valid views of the same encoder.

\paragraph{E2: unsupervised segmentation quality (C4).}
This experiment tests whether motion-word quantization cuts continuous streams into behavior-aligned proposals. Under a seeds pseudo-ground-truth protocol, boundary IoU against held-out annotations is \textbf{0.458 $\pm$ 0.049}, versus an equal-segment-count random-cut null of 0.323 $\pm$ 0.022 (Monte-Carlo expectation over 200 random cuts); boundary F1 is reported alongside (0.343 vs.\ the uniform-window 0.396), and the fixed-grid probe that is statistically on par with SMQ ($0.453 \pm 0.027$) is reported in Section~\ref{sec:ablation}.

Figure~\ref{fig:segqual} visualizes the resulting boundaries against the seed pseudo-ground-truth on the lowest- and highest-IoU episodes.

\begin{figure}[pos=tbp]
  \centering
  \includegraphics[width=\linewidth]{fig3_segmentation_qualitative.pdf}
  \caption{Qualitative segmentation boundaries on the two extreme episodes of the public-real tier (lowest- and highest-IoU of the four episodes). Upper band per episode: seed pseudo-ground-truth segments (orange; rule-engine seeds, confidence $\geq 0.8$ and duration $\geq 0.5$\,s---a weak reference, not human annotation); lower band: SMQ predictions (teal; end-to-end K=8, epoch-30 checkpoint). The lower panel aggregates matched IoU against the equal-segment-count random-cut null over all four episodes (0.4577 $\pm$ 0.0488, 4/4 episodes above baseline), consistent with the ablation in Section~\ref{sec:ablation}. Boundary matching is label-agnostic (Hungarian matching~\citep{kuhn1955hungarian} on IoU), so band colors encode segmentation quality, not class semantics.}\label{fig:segqual}
\end{figure}

\paragraph{E3: anchor-guided label expansion (C5).}
This experiment tests whether a small rule-engine seed set expands toward taxonomy coverage via iterated pseudo-labeling. Anchor-guided prototype clustering reaches purity \textbf{0.5339} against a 0.3306 random-assignment baseline ($1.61\times$) and a 0.4858 all-majority baseline ($+4.8$\,pp---the binding comparison). Injecting 30\% label noise into the seeds degrades purity by only 3.1\,pp (0.5025$\pm$0.0063, still $1.52\times$ baseline), and enlarging the seed ratio from 25\% to 100\% changes purity by less than 1\,pp at $K\in\{5,10\}$ but costs 2.4\,pp at $K{=}14$, the best cell of the $k$-means sensitivity sweep (the main run is the deterministic class-mean configuration; full sweep in the released evidence archive)---prototype quality saturates early in seed ratio at small $K$, which is what makes a small seed budget practical.

\paragraph{E4: semi-supervised self-training on the public-real tier (C6).}
This experiment tests whether self-training converts pseudo-labeled coverage into usable classification under scarce labels. Iterating from round~0 to round~1 raises pseudo-label pool precision from 0.5125 to \textbf{0.691 $\pm$ 0.013} (an $+17.9$\,pp endpoint gain; the paired first-to-last statistic over all iterated rounds is $\Delta$ $+10.69\pm3.28$\,pp---the two numbers measure different spans and are both released); the precision-drop stopping rule fired after two consecutive degraded rounds, so round~1 is the reported operating point (Holm-corrected $p=0.090$ over the three-test family, raw $p=0.030$---direction-consistent but not significant after family correction). Under the pre-registered partial-class protocol on the Animal Kingdom canine subset (4 of 12 classes pass the sample-size gate), a frozen-backbone classifier with a re-trained head reaches 44.90\% overall (seed 42; three-seed mean 41.5\%, sample std $\pm$5.9\,pp) ($1.80\times$ the 25\% uniform-random baseline). The per-class breakdown (watch 100\% / track 23.5\% / jump 0\% / stay 0\%) accompanies this number: the aggregate is dominated by class imbalance and interprets only in the four-class partial caliber (other public-real results carry their own class-space labels).

\paragraph{E7: end-to-end pipeline at small budget (C6).}
This experiment tests whether the complete pipeline, run end-to-end from a small annotation budget, approaches full supervision. The arms follow a corrected protocol (revision R16 of the repository's experiment ledger): the final classifier is trained on seed ground-truth labels plus pseudo-labels for the confidence-filtered pool clips---never on pool ground truth---and the iteration stops on convergence or the iteration budget (the precision-drop rule monitors pool precision against a rule-derived consensus reference, which this tier's label space does not provide; consuming true labels there would be oracle leakage---Section~\ref{ssec:proto}). On the 12-class-protocol public-real tier (9 classes with samples; 197 self-extracted clips; anchors and pool restricted to the training split, validation isolated), at 2 clips per class (18 of 141 training clips, $\approx$13\%; ten seeds) the pipeline reaches \textbf{9.8\% $\pm$ 7.5} top-1---at the tier's 11.1\% chance level---against 33.93\% under full supervision; the generic-SSL control reaches 15.2\% $\pm$ 5.1 and the scratch control 8.4\% $\pm$ 10.7; after per-tier Holm correction the only significant same-seed effects are on macro-F1, both favoring the generic-SSL arm (warm trailing by $5.1$\,pp at 2 clips/class, corrected $p=0.023$, and $6.6$\,pp at 4, corrected $p=0.049$); the warm-versus-scratch differences are not significant. The prototype-path consensus gate is inert on this taxonomy (no standing class; Section~\ref{ssec:proto}), so pool filtering reduces to the confidence margin alone. The iterated pseudo-labeling is the point of failure: pool pseudo-label precision, measured post hoc against training-split labels as a released diagnostic only, averages 0.11 across seeds, and a ground-truth-free per-round recalibration of the gate on anchor-side margins does not rescue it ($8.9$--$12.5\%$, three seeds). The negative boundary\footnote{An earlier internal protocol variant trained the final head on the pool clips' ground-truth labels and consumed oracle precision for stopping; its higher low-budget numbers were never published and are withdrawn. The correction, re-run evidence, and errata are documented in the released repository.}: on the canine tier, confidence-filtered pseudo-labels at a 13\% budget do not yield usable supervision---consistent with the label-alignment diagnosis of Section~\ref{sec:conclusion} (L9) and with the saturated-scorer mechanism of the active-learning study (E5, Section~\ref{sec:ablation}). The supervised references and comparisons on this tier (E4, the five-arm study in Section~\ref{sec:comparison}) are unaffected, and the cross-domain retention evidence (E9) holds under the same corrected protocol; Figure~\ref{fig:budget} collects the retention points across all exercised tiers in one view.

\paragraph{E7b: pre-registered replication on the expanded v2 tier.}
To test the data-bottleneck attribution instead of asserting it, we pre-registered a multi-segment expansion of the public-real tier (protocol PSD-AKV2-PREREG-001, frozen before the build; protocol and results landed in the same repository batch): the same 329 canine videos re-extracted as up to four contiguous segments each under a per-frame label-consistency gate ($\geq$0.80), yielding 352 clips (256 train / 96 val) in an eight-class space (bark segments fail the gate). All low-budget arms rerun under the corrected pseudo-label protocol of E7 at ten seeds. Leakage boundary: the warm-start backbone was supervised on v1 training-split videos, and v2 adds new segments of the same videos to its training pool; validation videos remain disjoint, so validation is clean, but any warm-arm advantage partly co-occurs with the warm arm's source-video exposure and is not fully attributable to the method alone. The frozen decision rule fires on raw top-1 ($+3.57$\,pp), but that gap is confounded by the class-space change (nominal chance rises from the 12-class $8.33\%$ to the 8-class $12.5\%$, a $+4.17$\,pp shift; the realized 9-class chance on v1 is $11.1\%$, E7); we therefore add a same-space control (post-hoc, not a pre-registered endpoint)---reweighting the v1 full-supervision arm's per-class accuracies over the eight classes v2 retains, on v1's validation clips, yields 25.96\% (released artifact \texttt{r12-holm-eightclass}, \texttt{eightclass\_control} block; errata for the artifact's remaining blocks are documented in the released repository); the control removes the chance-rate confound but not the clip-set difference---so the same-space rise is $+11.54$\,pp (above-chance margin $13.5\to25.0$\,pp; above-majority $-4.8\to+3.1$\,pp). Under this control the top-1 data are consistent with the data-bottleneck attribution, and macro-F1 moves the same way ($16.5\%$ to $19.3\%$), so both metrics agree. The raw endpoint is numerically sensitive: re-running the identical supervised reference with a different label encoding shifts it by 2.1\,pp ($37.50\to35.42$; two validation clips), so the raw margin above the $+3.0$\,pp line is within solver-path noise and the attribution rests on the same-space control, whose margin is an order of magnitude larger. The corrected low-budget arms do not replicate the earlier end-to-end reading: at 2 clips/class where a class has training clips (14 of 256, $\approx$5.5\%; the sit class has none), the warm-start arm reaches 13.1\% $\pm$ 5.3, the generic-SSL arm 17.6\% $\pm$ 4.3, and the scratch arm 13.8\% $\pm$ 5.6, against a 12.5\% chance rate---near chance, as on v1; at 4 clips/class (10.9\%, 28 of 256; the point plotted in Figure~\ref{fig:budget}) the warm arm reads 16.6\% $\pm$ 6.8, i.e.\ 44.2\% retention of the archived 37.50\% v2 full-budget reference, with generic-SSL at 14.6\%---ordered as at spc2 (spc: samples per class); the only significant same-seed effect in this family is macro-F1, where the warm arm beats scratch ($+4.6$\,pp, all ten pairs, per-tier Holm-corrected $p=0.012$). v1 remains the primary tier; v2 is reported alongside, never merged. As a diagnostic control (not part of the pre-registered endpoints), we also end-to-end fine-tune the backbone on v2: full-budget fine-tuning reaches 32.99\% $\pm$ 3.94 top-1 (three seeds)---\emph{below} the frozen-head reference on both metrics (top-1 37.50\% archived, 35.42\% under the re-encoding noted above; macro-F1 19.33 vs.\ 16.73); at 2 clips/class fine-tuning reaches 9.72\% $\pm$ 1.59, comparable to the near-chance pseudo-label arms. Both metrics agree that unfreezing does not help at this data scale---the full-budget frozen-versus-tuned comparison that the decoupled design prescribes.

\paragraph{E7c: skeleton-quality attribution (pre-registered).}
The negative boundary above leaves one attribution open: is the ceiling set by skeleton quality or by label alignment? We pre-registered the test (protocol PSD-SA-PREREG-001, frozen before the run): re-extract all 197 public-real videos with a substantially stronger estimator---YOLO11x-pose~\citep{jocher2024yolo11} fine-tuned on the dog-pose benchmark (up to 100 epochs, early-stopped at 89; pose mAP50 0.834 versus the production extractor's smaller model; extractor-training logs released with the K9 training repository)---and re-evaluate the full-supervision reference on identical clips, splits, and labels. The stronger skeletons do \textbf{not} raise the ceiling: 30.36\% versus 33.93\% ($-3.57$\,pp, below the pre-registered $+3$\,pp line; verdict LABEL\_BOTTLENECK; single-run references on both sides, no seed dispersion reported), so the low-budget arms are skipped per the frozen rule. Together with the alignment, backbone, and budget controls of Section~\ref{sec:ablation}, this triangulates the attribution from a fourth direction (single-run evidence on both sides; the $-3.57$\,pp shift lies within this tier's three-seed spread of $\pm$5.9\,pp and is read as decision-grade per the frozen rule, not as a precise effect estimate): on this tier the binding constraint is the video-level single-label annotation protocol (one label per clip over multi-behavior footage, misaligned with frame-accurate skeletons), not the skeleton extractor---a tier-specific diagnosis, not a claim that coarse granularity is universally prohibitive: PanAf500 reaches high retention under frame-level majority-voted clip labels precisely because its full-budget reference is separable (60.0\%, pool precision viable), whereas the canine tier's reference (33.93\%) with near-random pool precision (0.11) cannot absorb label noise of any granularity---the same pretext-learnability gradient that governs the E9 series (Section~\ref{sec:experiments}). What a deployment would iterate on is the label alignment layer, not re-extraction (the decoupled design's premise).

\paragraph{Warm-start small-budget usability (C6).}
On the synthetic-offset tier (pre-registered offset distribution, noise\_std$\,=\,$0.10; the offset diagnostic precedes the warm-start experiment), initializing the semantic layer from full-budget weights makes a 20-clip budget usable: \textbf{82.0\% $\pm$ 4.3} top-1 on 22 classes (three seeds), against 95.7\% for the same warm-start protocol at the 200-clip budget point (91\% of the 220-clip offset pool)---86\% of full-budget accuracy from 9\% of the labels---and against 41.2\% for the frozen, un-adapted backbone on the same offset validation set, so small-budget warm-start fine-tuning nearly doubles the un-adapted model ($1.99\times$) under distribution shift. The budget accounting must be read in full: the initializing weights come from a model trained to saturation on the \emph{clean} synthetic tier (2{,}200 clips of the same generator, no offset; zero-shot on the offset domain it reaches only 41.2\%), so the 20 offset-domain clips are the marginal cost on top of a fully labeled source domain, and the experiment quantifies adaptation under distribution shift---not learning a taxonomy from 20 clips alone. The cold-start comparison at the same 20-clip budget (8.0\% $\pm$ 3.0, no pretrained initialization) isolates exactly this marginal contribution.

\paragraph{E6: taxonomy-transition cost (C1).}
This experiment tests whether absorbing a taxonomy change under decoupling costs less than full-pipeline retraining at matched accuracy. On the synthetic tier, under a pre-registered synthetic $\mathcal{Y}\to\mathcal{Y}'$ remap (released with the code), the baseline arm retrains the full ST-GCN pipeline while the decoupled arm freezes the $\mathcal{Y}$-era backbone and rebuilds only the semantic head---three seeds per budget (small: $n{=}30$ clips, CPU; full: $n{=}100$ clips, GPU), identical annotation units across arms. The decoupled semantic-layer rebuild averages 31.1\,s of wall-clock retraining versus 188.7\,s for full-pipeline retraining ($6.07\times$; same-seed paired ratios never below $3.99\times$, seeds 42--44); convergence epochs improve by a further $2.18\times$; and accuracy differences fall inside the pre-registered noise band ($-0.91$\,pp full / $+2.27$\,pp small, both $<2.3$\,pp). Annotation units are identical across arms by construction, so the saving is computational rather than labeling-driven. Following our pre-registered conservative-bound rule we claim \textbf{$\geq 3\times$ lower retraining cost}, with the accuracy equivalence stated as statistical equivalence rather than any directional advantage. The absolute wall-clock values are benchmark-specific; the claim concerns the ratio, whose practical value scales with the size of the pipeline being retrained. The attribution boundary: the wall-clock ratio reflects the frozen-versus-retrained architecture distinction and is available to any design that freezes a backbone and retrains its head---what the decoupled design contributes at a transition is not the ratio itself but the frozen path's accuracy and label economics: the semantic layer rebuilds from rule-engine seeds plus unsupervised proposals rather than full re-annotation of $\mathcal{Y}'$ (the E6 protocol here maps the same labeled units to both arms, so the seed-budget advantage of $\Omega$ is not isolated by this experiment and is exercised instead by the small-budget arms of Sections~\ref{sec:experiments}).

\paragraph{E6-real: taxonomy-transition cost on the human benchmark (pre-registered replication).}
A pre-registered replication (protocol PSD-NTU-TRANS-001, frozen before any run) transfers this two-arm design to NTU60 under a frozen ten-merge taxonomy transition ($60\to49$ classes, the coarse-grained reporting scenario), with the identical joint-level MLP backbone in both arms (E9b/c convention; decoupled = frozen Y-era pretext plus a linear head retrained on $\mathcal{Y}'$ labels; coupled = the same architecture retrained end-to-end from scratch on $\mathcal{Y}'$ labels; seeds 42--44). The pre-registered cost endpoint \textbf{fails}: the decoupled arm's head retraining (the E9-series CPU logistic-regression solver, 466\,s median) is not cheaper than end-to-end retraining of the small coupled MLP (35\,s)---a $13\times$ cost asymmetry at full budget ($2.3\times$ at 10\% labels) whose substantive reading is scale: when the entire coupled pipeline retrains in 35\,s, there is little retraining cost for decoupling to save, exactly the scaling caveat above; the frozen protocol pinned architecture, epochs, budgets, and seeds but not solver families, and we do not re-run with a matched solver post hoc. The pre-registered accuracy expectation is \emph{not} confirmed as stated: the protocol anticipated the coupled arm may dominate at full supervision (per the L10 gradient) while the decoupled arm may dominate at 10\% labels; the decoupled arm was in fact more accurate at every budget ($+3.9$\,pp at full, $+11.3$\,pp at 10\% labels, both outside the $\pm2.3$\,pp equivalence band)---the anticipated coupled-arm dominance did not materialize, with the solver-family confound (disclosed above) applying to both endpoints. On real data the measured value of the decoupled pattern at this scale is accuracy, not wall-clock, and per the frozen decision rule the $\geq3\times$ claim remains scoped to the synthetic tier with its matched-solver arms.

\begin{figure}[pos=tbp]
  \centering
  \includegraphics[width=\linewidth]{fig5_budget_retention.pdf}
  \caption{Budget--retention summary: each point is one tier's PSD pipeline at a given fraction of that tier's full-labeled pool, plotted as the fraction of its own full-budget top-1 retained (error bars: seed std; corrected pseudo-label protocol; the NTU60 bar, $\pm$0.3\,pp, is smaller than its marker). Tier groups are never connected or pooled; where budgets nearly coincide, points are displaced along the log axis for legibility and the displacement is listed per point (v1 12.8\% shown at 13.9\%, v2 spc4 10.9\% shown at 11.5\%, synthetic-offset 9.1\% shown at 8.0\%, UCF101 10\% shown at 8.8\%, PanAf500 10\% shown at 13.2\%, NTU120 10\% shown at 16.5\%); all plotted retention values, error bars, and the ordinal spacing of the retention axis are unaffected by the displacement. The pattern is tier-dependent: retention is high on the human benchmarks (NTU60 90.7\%, NTU120 88.9\%), on the animal public benchmark PanAf500 (88.7\%), and on the synthetic-offset tier, but collapses toward chance on the canine public-real tier at 13\%/5.5\% budgets (E7/E7b) and drops to 66.6\% on the independent UCF101 benchmark, whose full-budget reference is itself weak (23.1\%)---the low-resource behavior of iterated pseudo-labeling is bounded by pseudo-label quality, not by budget arithmetic.}\label{fig:budget}
\end{figure}

\subsection{Comparison studies}\label{sec:comparison}

\paragraph{Implementation-equivalence verification on NTU60.}
Because AimCLR, SMQ, and the self-training machinery are re-implemented in this repository rather than imported, we verify implementation correctness on the human-domain benchmark (NTU RGB+D 60~\citep{shahroudy2016ntu}) before applying the pipeline to animal data. Following the linear-evaluation protocol of AimCLR~\citep{guo2022aimclr} on NTU60 cross-subject, each stream (joint, bone, motion) is pretrained for 300 epochs and evaluated with a frozen encoder and a fully-connected classifier. The three streams are then fused with the official fixed weights (joint 0.6, bone 0.6, motion 0.4). The re-implementation reaches 74.30\% (joint), 71.51\% (bone), and 67.84\% (motion) top-1, and the three-stream fusion attains \textbf{77.97\% top-1 (95.78\% top-5; $n=16{,}487$)}---exceeding our pre-registered acceptance line of 77.18\% and lying within 0.93\,pp of the 78.9\% reported in the AimCLR paper body. We therefore treat the adapted implementation as equivalent to the official one for the purposes of this study.

\paragraph{E9: low-resource retention on the human benchmark (C6, pre-registered).}
The animal-domain low-resource claim would be weak evidence if it did not generalize, so we pre-registered (protocol PSD-NTU-PREREG-001, frozen before the experiment) a budget-shrinkage test on NTU RGB+D 60 cross-subject~\citep{shahroudy2016ntu}. Holding the frozen epoch-300 joint pretext backbone fixed, a 10\% class-stratified subset of the 40{,}091 training clips in the provider's frame-50 export (counts differ slightly from the original release) yields 66.05\% with a linear head alone, and \textbf{67.53\% $\pm$ 0.24 (ten seeds, after a 3-seed initial readout was expanded per the E9-series convention; both readouts are released) once the PSD semantic pipeline (anchor-guided clustering with iterated confidence-filtered pseudo-labeling) is applied under the corrected protocol of E7---90.7\% of the full-budget linear-probe reference of 74.45\%}, clearing the pre-registered generalization line of 90\% (the denominator is arm (c)'s linear probe, fixed at protocol freeze; under the three-stream fusion caliber of Section~\ref{sec:comparison} the same arms read 86.6\%)---though the linear-only arm alone already retains 88.7\% of the reference, so the pipeline's margin over a trivial budget-restricted head is modest: the pseudo-label iteration contributes $+1.48$\,pp over the linear-only arm (every seed arm above it; one-sample $t$ (raw $p<10^{-7}$; Holm-corrected $p<10^{-6}$), Wilcoxon (raw $p=0.002$; corrected $p=0.012$) under the series-wide eight-test family defined in Section~\ref{sec:implementation}), on a labeled subset $28\times$ larger than the animal-tier training pool (4{,}009 vs.\ 141 clips). An external semi-supervised baseline (protocol PSD-SSL-BASE-001, frozen before the run) bounds this margin from below in the same harness---same frozen features, same 10\% subset, same evaluation: a Mean-Teacher-style~\citep{tarvainen2017mean} pool-distillation arm (soft pseudo-labels, $\tau{=}0.95$ confidence mask, five teacher--student rounds, ten seeds) reaches \textbf{46.79\% $\pm$ 0.61}, so the PSD pipeline exceeds the classical external SSL baseline by $+20.7$\,pp (10/10 paired seed wins, Wilcoxon raw $p=0.002$; Holm-corrected $p=0.012$ in the same family), while the baseline's confidence mask admits only 15{,}758 of 36{,}082 pool clips---consistent with the mechanism reading above, since near-full pool restoration is exactly what the confidence-thresholded baseline cannot do. On the weak-reference domain (UCF101) the same baseline lands at parity with its teacher (14.17\% vs.\ 14.04\%; PSD $+1.2$\,pp, n.s.), so the external-baseline disadvantage tracks the reference-quality gradient across both domains. Retention ratios are protocol-dependent---a frozen-probe budget-accuracy curve is flatter by construction than a fine-tuned pipeline's---so this arm supports the pipeline's own budget behavior and is not a comparison against any published method. Four disclosures bound the reading: the confidence-filtered pool adds $\approx$35.4k pseudo-labeled clips per seed---the near-restored training-set size \emph{is} the mechanism---the gate accepts $\approx$98\% of the unlabeled pool at this scale (and the prototype-path consensus gate is inert on the numeric NTU taxonomy), so this arm cannot separate filter quality from pool-size restoration (Section~\ref{sec:conclusion}, L11); the iteration's precision-drop stopping rule is disabled on this tier (no rule-derived consensus reference exists), and pool precision against the official training labels ($\approx$0.69) is reported only as a post-hoc diagnostic that never enters the loop's control path; the 10\% subset is a single stratified draw (seed 42), so the $\pm$0.24 spread covers self-training seeds but not subset resampling---a post-hoc robustness diagnostic (not a pre-registered endpoint) that redraws the subset under four alternative seeds places the linear arm's subset-resampling spread at 65.62\% $\pm$ 0.48 (original draw 66.05\%, within one resample standard deviation; the decoupled arm's subset sensitivity is not measured---each resample would require a full pipeline re-run); and the final head uses StandardScaler plus logistic regression (tol $10^{-3}$) with the self-training head on GPU---scale adaptations, all other hyperparameters identical to E7.

\paragraph{E9b/E9c: cross-domain retention gradient (pre-registered).}
To test whether the retention behavior above is a property of the pipeline rather than of one benchmark, we pre-registered two further budget-shrinkage replications (protocols PSD-NTU120-PREREG-001 and PSD-UCF101-PREREG-001, frozen before any run). On NTU RGB+D 120~\citep{liu2020ntu}, the official 3D-skeleton application was denied, so the replication uses the PYSKL-release HRNet 2D joint annotations (same videos and labels, different extractor); the pretext is a joint-level MLP over all joint coordinates (80 epochs, frozen 256-d penultimate)---a disclosed deviation from the frozen ST-GCN plan, adopted so one pretext architecture runs unchanged across both 17-joint benchmarks, with backbone quality absorbed by the full-budget reference arm as before. At a 10\% stratified budget (ten seeds, after a 3-seed initial readout was expanded per the E9-series convention---the ten-seed convention was adopted 2026-09-07 as a dated amendment, both readouts released; as in E9, spreads cover self-training seeds but not subset resampling), the pipeline reaches \textbf{48.3\% $\pm$ 0.8} top-1 against the 10\%-linear arm's 45.4\% ($+2.9$\,pp; every seed arm above the linear arm; one-sample $t$ vs.\ the deterministic linear arm, $t{=}12.0$, raw $p<10^{-5}$, Wilcoxon raw $p=0.002$; Holm-corrected over the series-wide eight-test family (Section~\ref{sec:implementation}; three tiers$\,\times\,$two statistics plus the two external-SSL tests)---corrected $p<10^{-5}$ and $0.012$; conclusions are invariant to any narrower family choice) and the full-budget reference's 54.3\%---\textbf{88.9\% retention, in the pre-registered PARTIAL band (85--90\%)}. On UCF101~\citep{soomro2012ucf101}---a genuinely independent third domain (YouTube-sourced video, 101 classes, HRNet 2D skeletons~\citep{sun2019hrnet})---the same protocol yields \textbf{15.4\% $\pm$ 1.9} against a 14.0\% linear arm ($+1.4$\,pp; 8/10 seed arms above the linear arm; raw $t$ $p=0.049$ and Wilcoxon raw $p=0.053$, Holm-corrected $p=0.141$ for both over the series-wide family---neither significant; direction-consistent, not statistically significant) and a 23.1\% reference: \textbf{66.6\% retention, below the 85\% line, reported as a negative boundary}. The external SSL baseline (Amendment 1 of PSD-SSL-BASE-001) extends to this domain: the Mean-Teacher arm reaches \textbf{14.17\% $\pm$ 0.87}---parity with its own teacher, the linear head (14.04\%)---so the PSD pipeline's advantage over the external baseline is $+1.2$\,pp (8/10 wins, Wilcoxon raw $p=0.047$, Holm-corrected $p=0.141$; not significant): where the full-budget reference is weak, confidence-thresholded pool distillation adds nothing, the same pretext-learnability gradient now observed on the external-baseline axis. Read together, the three cross-domain points form a monotone gradient: retention 90.7\% $\to$ 88.9\% $\to$ 66.6\% tracks the full-budget reference quality 74.45\% $\to$ 54.3\% $\to$ 23.1\% (all E9-protocol linear-probe references; the AimCLR official-protocol re-implementation of the NTU60 reference, with its FC-softmax head, is 74.30\%, Section~\ref{sec:comparison}), indicating that what the frozen pretext learned about class-separable structure---not the budget arithmetic---bounds how far pseudo-label self-training can be stretched. The gradient is stated in falsifiable form: retention is monotone in full-budget linear-probe reference quality, with the high band (85--90\%) holding where the reference is separable and collapse where it is not; any tier whose full-budget reference is strongly separable ($\geq$70\%) yet retains $<85\%$ at a 10\% budget would break it---the pre-registered working-dog pilot is designed to supply such a test in the target domain. This is the same mechanism boundary diagnosed on the canine tier (E7, Section~\ref{ssec:proto}; the alignment analysis of Section~\ref{sec:conclusion}, L9), now observed on two additional public benchmarks; the low-resource claim is therefore scoped to tiers whose pretext features support linear separability at full budget, and the retention line is not claimed as a domain-universal property.

\paragraph{E9d: first animal-domain public benchmark (PanAf500, pre-registered).}
The mechanism boundary above predicts that retention depends on what the pretext learned, not on taxonomy---so it should extend beyond the human domain wherever a public animal benchmark exists. PanAf500 (great-ape natural behavior, in-the-wild video, 500 clips with official train/val/test splits and nine behavior classes) is the public animal benchmark we could identify, and we pre-registered its budget-shrinkage test (protocol PSD-PANAF-PREREG-001, frozen before any run)---\emph{after} the E9b/c gradient was observable but \emph{before} the PanAf run executed. Its full-budget reference (60.0\%) lies below the $\geq$70\% falsifier threshold of the gradient's strong form, so this run is a directional consistency check rather than a potential falsifier: the observed 88.7\% retention (PARTIAL band) is consistent with the high band predicted for separable references. Skeletons are model-extracted: a YOLO11x-pose model fine-tuned on the ape subset of AP-10K (chimpanzee/gorilla/orangutan, 17-keypoint quadruped topology mapped onto the 24-slot layout with seven dead slots ($24-17$), subsuming the canine tier's four) infers keypoints on bbox-cropped video frames (released quality: 82.1\% mean frame detection rate, 0.59 mean confidence; extraction noise enters all arms equally). The pretext, arms, budget, and decision rule mirror E9b/E9c exactly. At a 10\% stratified budget (ten seeds), the pipeline reaches \textbf{53.2\% $\pm$ 9.1} against the 10\%-linear arm's 52.0\% and the full-budget reference's 60.0\%---\textbf{88.7\% retention, in the pre-registered PARTIAL band} (a 3-seed early readout of 93.3\% fell into the CONFIRMS band and is superseded by the 10-seed number under the frozen rule; both are released). The test split (75 clips, held out of the decision rule) yields 56\% as a non-decision secondary number. Two disclosures bound the reading: the nine-class space is long-tailed (three classes have $\le$7 clips; five classes contribute a single seed clip at the 10\% budget), which inflates seed variance; and behavior labels are frame-level annotations majority-voted to clip level---the same coarse label granularity diagnosed on the canine tier. PanAf500 completes the cross-domain picture at four points spanning two kingdoms: retention stays in the high band wherever the full-budget reference is linearly separable (NTU60 90.7\%, NTU120 88.9\%, PanAf500 88.7\%, whose reference quality sits between the two NTU tiers) and collapses where it is not (UCF101 66.6\%, reference 23.1\%)---budget retention is a property of pretext learnability, not of species or domain. The PanAf pipeline-vs-linear gap is not carried by the evidence at any tested power: at ten seeds it reads $+1.2$\,pp with 5/10 seed wins, and a post-hoc seed expansion to thirty seeds (Amendment 2 of PSD-PANAF-PREREG-001, frozen before the run; seeds 42--71, both readouts released) gives 51.5\% $\pm$ 8.6 against the same linear arm---13/30 wins, 3 ties, Wilcoxon $p=0.54$, retention 85.8\%, still in the PARTIAL band with the verdict unchanged---so the gap is reported as directional only, and even $n{=}30$ does not separate the arms on this long-tailed nine-class space; as on NTU60, PanAf's retention is carried predominantly by the frozen-feature linear head, and no pipeline advantage is claimed on this tier.


\paragraph{E6-real-2: backbone-scale replication with matched solvers (pre-registered).}
TRANS-001's coupled arm retrains in $\approx$35\,s, so its cost endpoint was arithmetically unpassable---the replication could not have confirmed the cost claim at that scale. A second pre-registered replication (protocol PSD-NTU-TRANS-002, frozen before the run; same $\mathcal{Y}\to\mathcal{Y}'$ merge, 60$\to$49 classes) therefore scales the backbone to the full ST-GCN of the equivalence verification and adds a third arm that decouples the two confounds TRANS-001 carried: the matched-solver arm (M) retrains a linear head on the coupled arm's own penultimate features with the decoupled arm's solver, isolating representation quality from retraining cost. Full budget (40,091 clips, seeds 42--44): the coupled arm retrains in a median of 9897\,s against the decoupled head's 626\,s (15.8$\times$; verdict \textbf{PARTIAL} per the frozen rule), accuracy 79.8\% (coupled) vs 75.3\% (decoupled, gap -4.5\,pp). A first run at the nominal scaled learning rate collapsed to a constant predictor mid-training on all three seeds (NaN) and was discarded under a NaN watchdog added to the protocol; the reported run uses half the linear-scaled learning rate with the same schedule. All three coupled-arm seeds include host-restart segments (host restarts under consumer-host RAM pressure); wall-clock is the sum of contributing segments per seed, with the per-segment decomposition in the released artifact. The verdict is robust to the accounting: a final-segment-only bound (6.2$\times$ still clears the $3\times$ line. At the 10\% budget the same replication reads 32.5$\times$ (coupled 1520\,s vs decoupled 47\,s; ratio computed on unrounded seconds) with a +2.6\,pp decoupled advantage (PARTIAL at the 10\% tier).

\paragraph{Post-hoc robustness expansion: ten-seed replication of the matched-solver comparison (protocol PSD-NTU-TRANS-002-MEXP, dated).}
Per a dated addendum, frozen before the expansion was run, the matched-solver comparison was replicated over ten seeds (42--51; the three seeds of the preceding paragraph are reused unchanged). At the frozen three-seed budget the matched-solver arm already outscored the coupled arm (mean $82.6\%$ vs $79.8\%$). Over the ten seeds it exceeds the coupled arm on every seed (mean $82.6\%$ vs $79.7\%$; $+2.9$\,pp; paired Wilcoxon signed-rank $p=0.002$; sign test $10/10$, $p=0.002$). This upgrades the mechanism reading from a three-seed observation to a paired test: end-to-end training under-exploits its own representation---a linear head retrained on the coupled arm's frozen penultimate features is more accurate than the coupled arm's end-to-end-trained output. Seeds with host-restart segments carry a per-segment wall decomposition in the artifact; wall-times do not enter the frozen cost verdict. The frozen cost-ratio verdict is unaffected by the expansion.

\paragraph{Segmentation strategy.}
A three-arm ablation under the seeds pseudo-ground-truth protocol (Section~\ref{sec:ablation}) shows motion-word quantization (0.458$\pm$0.049) ahead of an equal-segment-count uniform-window arm (0.399$\pm$0.035) and the random-cut null (0.323$\pm$0.022), passing all three pre-registered gates; a fixed-grid probe (0.453$\pm$0.027) is statistically on par with SMQ.

\paragraph{End-to-end comparison against training from scratch.}
To complete the comparison matrix, we ran five new arms on the 12-class public-real protocol (9 classes with samples; 197 self-extracted clips; three seeds; identical budgets), reported alongside the two frozen-backbone conditions at full training-split budget. The seven reported conditions are: the two frozen-backbone head-retrain arms (warm vs.\ random backbone, trained on the full training split), the same two initializations with \emph{full} fine-tuning (no freezing), and three trivial baselines. Under the frozen protocol---the usage mode our decoupled design prescribes---the pretrained backbone outperforms a random one by $+8.3$\,pp (35.71\% vs.\ 27.38\%; seed-level std is reported as dispersion, not as a confidence interval). Under full fine-tuning, however, the advantage disappears (36.90\% vs.\ 36.90\%; the from-scratch arm even attains the best macro-F1, 0.119 vs.\ 0.103---macro-F1 over the 12-slot label space with absent classes counted as zero, same caliber across arms), which corroborates the gradient-ablation crossover finding of Section~\ref{sec:ablation}: past the crossover annotation budget, the pretrained prior adds no value \emph{as a fine-tuning initialization}---the frozen-head usage mode the decoupled design prescribes retains the prior's benefit at full budget (the E6-real replication, Section~\ref{sec:experiments}). Most tellingly, every deep arm and a flatten-plus-logistic-regression baseline cluster in a narrow 27--37\% band (random 6.0\%, majority 25.0\%): the ceiling of this tier is set by the self-extraction data bottleneck (197 clips, single-label, severe imbalance), not by the choice of model. The study is therefore evidence of a data bottleneck and of where pretrained priors do and do not help, not a performance claim. Reference calibers on this tier differ by protocol and are not interchangeable: 33.93\% (the unified-label 12-class reference of E7), 35.71\% and 36.90\% (this five-arm study), and 37.50\% and 35.42\% (the v2 re-extraction under two label encodings) span up to 3.6\,pp---a caliber effect, not a model effect; within-protocol arm comparisons are unaffected.

\paragraph{Scope statements.}
The NTU60 equivalence verification covers the AimCLR pretraining and three-stream fusion protocol; SMQ is validated through pre-registered ablation comparisons rather than against an official baseline protocol (none is published), and the TCL-style self-training is delivered as an equivalent iteration loop since no official implementation exists (Section~\ref{ssec:selftrain}). Two further registered comparisons remain outside the present evaluation and are not claimed anywhere in this paper: the AimCLR++ backbone variant (80.9\% NTU cross-subject per the official repository's linear-evaluation table) as an alternative physics-layer encoder, and exhaustive sweeps over the confidence threshold $\tau$ and cluster count $K$.

\begin{table}[pos=t]
\centering
\caption{Main results under the three-tier protocol (plus the public human-domain benchmark used for implementation-equivalence verification). Tiers are never mixed; the real-K9 row is blank because that tier is not exercised in this study (Section~\ref{sec:datasets}). Per-class breakdowns accompany the four-class public-real aggregate (Section~\ref{sec:conclusion}).}
\label{tab:main}
\small
\hyphenpenalty=10000
\begin{tabularx}{\linewidth}{>{\raggedright\arraybackslash}p{2.7cm}>{\raggedright\arraybackslash}X>{\raggedright\arraybackslash}p{1.2cm}>{\raggedright\arraybackslash}p{4.1cm}}
\hline
Tier & Task / protocol & Metric & Result \\
\hline
Synthetic & 22-class, full budget (matched 50-ep) & top-1 & 96.6\% (single run, seed 42) \\
Synthetic-offset & 22-class, 20-clip warm-start & top-1 & 82.0\% ($\pm$4.3) \\
Synthetic & taxonomy transition, decoupled vs.\ full & wall-clock & 31.1\,s vs.\ 188.7\,s ($6.07\times$; $\geq 3\times$ claimed) \\
\hline
Public-real & $k$-NN subject probe (representation) & top-1 & 20.89\% ($2.51\times$ random) \\
Public-real & segmentation IoU (seeds pseudo-GT) & IoU & 0.458 $\pm$ 0.049 \\
Public-real & prototype purity (InterPet4D 7-label seed space) & purity & 0.5339 ($1.61\times$ random; $+4.8$\,pp majority) \\
Public-real & pseudo-label pool precision (round-1 endpoint) & prec. & 0.691 $\pm$ 0.013 (paired first-to-last n.s.) \\
Public-real & 4-class partial, frozen backbone + head & top-1 & 44.90\% (seed 42, best of 3; mean 41.5\,$\pm$\,5.9; per-class in Section~\ref{sec:conclusion}, L7) \\
Public-real & 12-class two-arm (warm vs.\ random backbone) & top-1 & 35.71\% vs.\ 27.38\% \\
Public-real & 12-class two-arm under full fine-tuning (null; five-arm study) & top-1 & 36.90\% vs.\ 36.90\% \\
Public-real & end-to-end PSD pipeline @ 13\% budget, corrected pseudo-label protocol (10 seeds) & top-1 & 9.8\% (chance 11.1\%; near-chance boundary, E7) \\
Public-real (v2 replication) & same corrected protocol, expanded tier @ 2 clips/class (10 seeds) & top-1 & 13.1\% (chance 12.5\%; E7b) \\
\hline
Public human benchmark & NTU60 xsub 3-stream fusion (equivalence) & top-1 & 77.97\% (line: 77.18\%) \\
Public human benchmark & PSD pipeline @ 10\% labels, corrected protocol (E9, retention) & top-1 & 67.5\% (90.7\% of linear-probe ref 74.45\%; line: 90\%, ten seeds) \\
Public human benchmark (HRNet 2D) & NTU120 replication @ 10\% labels, corrected protocol (E9b, retention) & top-1 & 48.3\% (88.9\% of probe ref 54.3\%; PARTIAL band 85--90\%, ten seeds) \\
Animal public benchmark (PanAf500) & PSD pipeline @ 10\% labels, corrected protocol (E9d, retention) & top-1 & 53.2\% (88.7\% of probe ref 60.0\%; PARTIAL band, ten seeds) \\
Independent benchmark (HRNet 2D, YouTube) & UCF101 @ 10\% labels, corrected protocol (E9c, retention) & top-1 & 15.4\% (66.6\% of probe ref 23.1\%; negative boundary, ten seeds) \\
\hline
Real-K9 & --- & --- & not exercised in this study \\
\hline
\end{tabularx}
\end{table}

